\section{立题依据及价值}

\subsection{研究背景与国内外研究现状}

大语言模型（Large Language Models, LLMs）已从单纯的文本生成系统演进为能够在现实世界中执行操作的自主智能体。现代LLM智能体可与外部工具交互、访问邮件与日历、查询数据库、浏览网页并执行代码\cite{schick2023toolformer, yao2022react}，这一演进催生了个人助理管理日程与通信、企业自动化系统处理文档与协调工作流、客服代理访问后端系统解决查询等诸多应用\cite{xi2024rise}。然而，能力提升的同时带来了显著的安全风险：当LLM不仅能生成文本，还能执行操作时，对抗性攻击的后果变得更加严重。能够发送邮件、转移资金或修改文件的智能体，与仅生成文本的聊天机器人相比，呈现出根本不同的威胁面。理解和评估这些漏洞对于LLM智能体的安全部署至关重要。

提示注入攻击（Prompt Injection）利用了LLM处理自然语言的基本特性。首先，指令与数据之间没有明确的边界\cite{greshake2023not, perez2022ignore}；其次，模型必须根据上下文理解来决定遵循哪些指令，而这一上下文敏感的指令遵循能力恰恰是智能体有效运作的关键。然而，同一能力也使智能体容易被操纵。攻击者将恶意指令嵌入到模型处理的内容中，由于模型无法仅根据来源或上下文可靠区分合法指令与对抗性指令，可能执行攻击者的命令。间接提示注入对LLM智能体构成了尤为严重的威胁\cite{greshake2023not}，攻击者不直接与智能体交互，而是将恶意指令嵌入智能体在运行过程中检索的外部内容（如邮件、文档或网页），导致数据泄露、未授权操作和输出完整性受损等后果。

当前，关于LLM自动化攻击的研究主要集中在越狱（Jailbreaking）方面，即诱使已对齐模型生成被禁止内容。GCG\cite{zou2023universal}、AutoDAN\cite{liu2024autodan}、PAIR\cite{chao2023jailbreaking}和TAP\cite{mehrotra2024tree}等方法已在越狱场景中被开发和广泛评估，在标准基准上实现了高成功率。相比之下，提示注入攻击在对抗性机器学习社区中受到的关注远少于越狱攻击。在提示注入文献内部，研究大多聚焦于直接提示注入场景与简单评估设置，多数工作依赖手工攻击或基本启发式策略\cite{liu2024formalizing, yi2025benchmarking}，而非自动化优化方法。评估场景通常限于简化任务（如带有恶意指令的文本摘要），在具有工具调用能力的智能体操作有状态环境的现实场景中，系统性的自动化攻击评估仍然匮乏\cite{debenedetti2024agentdojo}。

\subsubsection{基于越狱的LLM对抗攻击技术}

越狱攻击旨在诱使已安全对齐的LLM生成被禁止内容，根据模型访问级别可分为白盒和黑盒两大类。

\textbf{白盒优化方法。}基于梯度的攻击利用对模型参数的白盒访问，通过直接微分优化对抗输入。Zou等\cite{zou2023universal}提出GCG（Greedy Coordinate Gradient）算法，其核心思想是计算损失函数关于one-hot token嵌入的梯度，在每个对抗位置选取梯度幅度最大的$k$个候选token，通过批量前向传播评估每个候选替换的损失，选择使损失最小化的替换。该方法在白盒Vicuna模型上实现88\%攻击成功率，且优化的对抗后缀对黑盒系统表现出显著可迁移性（GPT-3.5达87\%、GPT-4达47\%），揭示了LLM对抗性脆弱的跨架构普遍性。Geisler等\cite{geisler2025attacking}引入投影梯度下降（PGD）方法攻击RAG系统，通过连续松弛重构离散token优化问题，在检索增强的生成管道中实现10倍计算加速下的可比成功率，并将攻击扩展到多文档检索场景。Wen等\cite{wen2023hard}建立梯度引导离散token搜索的通用范式，形式化离散优化中的硬约束问题，为后续GCG等算法提供了理论框架与算法基础。Liu等\cite{liu2024autodan}提出AutoDAN，采用遗传算法与分层交叉算子生成语义连贯、可解释的对抗提示。不同于GCG产生不可读的token序列，AutoDAN通过在语义空间中搜索，在保持与GCG竞争性成功率的同时生成更低困惑度、更难被检测的提示。Andriushchenko等\cite{andriushchenko2025jailbreaking}证明仅使用logprob反馈（无需梯度）的简单自适应攻击即可在GPT-4、Claude和Gemini等领先安全对齐模型上实现近100\%成功率，揭示当前对齐方法的脆弱性，表明现有安全训练未能建立真正的安全边界。

\textbf{黑盒优化方法。}Mehrotra等\cite{mehrotra2024tree}提出TAP（Tree of Attacks with Pruning）算法，采用三模型架构：攻击者LLM负责根据反馈迭代生成越狱候选，评估器LLM判断目标模型响应是否包含有害内容并评分，目标LLM作为被攻击对象。TAP通过树结构搜索（宽度$w$、深度$d$）探索多种攻击策略，利用on-topic分类器剪枝偏离目标方向的分支，在GPT-4上以平均29次查询实现80--94\%成功率。Sitawarin等\cite{sitawarin2024pal}提出PAL（Proxy-Guided Attack），利用代理模型（如Llama-2）上的GCG优化攻击后迁移至目标模型，结合智能重排序和多候选策略提高迁移成功率，首次系统研究了越狱攻击的代理辅助优化范式。Chen等\cite{chen2024instructzero}提出InstructZero，通过查询高效探索系统提示空间，在贝叶斯优化框架下，将黑盒越狱转化为连续提示嵌入空间上的优化问题，在无需梯度访问下实现与白盒方法竞争的性能。

\subsubsection{基于提示注入的LLM智能体攻击技术}

\textbf{白盒优化方法。}Pasquini等\cite{pasquini2024neural}提出Neural Exec，针对RAG系统中的延迟提示注入（deferred prompt injection），通过优化在检索条件下激活的"执行触发器"，不同于即时攻击，该触发器仅在特定检索上下文中激活恶意指令，在RAG系统上实现87\%攻击成功率，揭示了检索增强机制引入的新攻击面。Fu等\cite{fu2024imprompter}引入Imprompter，将GCG优化应用于生成导致智能体误用工具的混淆提示，首次将GCG从越狱场景适配到智能体工具误用场景，通过优化token序列使智能体调用非预期工具或传递错误参数。Zhan等\cite{zhan2025adaptive}系统研究自适应攻击对智能体防御的威胁，通过修改GCG和AutoDAN显式考虑防御机制（如结构化分隔符和指令层级），在StruQ、SecAlign等多种防御范式下均实现超50\%攻击成功率，证明当前防御在自适应攻击者面前的脆弱性。Wu等\cite{wu2025dissecting}将对抗鲁棒性评估扩展到多模态LM智能体，在视觉-语言模型智能体上实现跨模态攻击，揭示智能体攻击面在多模态设置下显著扩大，恶意图像可与文本注入协同增强攻击效果。

\textbf{黑盒优化方法。}Liu等\cite{liu2024prompt}提出HOUYI框架，将提示注入攻击系统化分解为三个组件：框架（framework，建立攻击上下文）、分隔符（separator，界定指令边界）和破坏器（disruptor，执行恶意指令），通过对这三个维度进行组合优化，在跨多个LLM应用（包括GPT-4和Claude）上实现高成功率。Zhang等\cite{zhang2024goal}提出G2PIA（Goal-guided Prompt Injection Attack），结合理论攻击原则（如权威伪装、紧急性框架等心理学原则）与LLM引导生成来产生多样化有效提示注入，通过自动化策略选择避免手工设计注入模板。Wang等\cite{wang2025agentvigil}提出基于蒙特卡洛树搜索（MCTS）的AgentVigil，将提示注入攻击建模为序列决策问题，以选择动作类型和注入内容作为搜索节点，在AgentDojo基准上实现71\%攻击成功率，代表当前智能体攻击的先进水平，但依赖手工设计的动作空间。Yu等\cite{yu2025promptfuzz}提出PromptFuzz，采用模板生成与LLM引导变异的两阶段模糊测试框架：第一阶段基于规则生成初始注入集合，第二阶段利用LLM对成功注入进行语义变异和增强，在多个智能体框架上实现攻击。

\textbf{其他攻击方法。}Labunets等\cite{labunets2025computing}利用微调API从闭源模型提取梯度等价信息进行优化攻击，绕过模型权重不可访问的限制，通过观察微调前后模型输出的差异反推梯度方向，在API-only设置下实现接近白盒的攻击效果。Shi等\cite{shi2025prompt}针对智能体推理循环中的工具选择决策实施两阶段优化攻击：第一阶段优化注入使智能体偏离原始推理路径，第二阶段优化注入使智能体执行攻击者指定的工具调用，揭示了智能体推理链的脆弱性。

\subsubsection{LLM智能体防御与评估技术}

\textbf{基于微调的防御。}Chen等\cite{chen2025struq}提出StruQ，通过结构化分隔符（如特殊token标记系统指令、用户输入和数据区域的边界）强制指令-数据分离，并在该分隔符格式上微调模型仅遵循系统指令区域的内容，在手工攻击下实现低于2\%攻击成功率，但在自适应GCG攻击下成功率仍可超50\%。Wallace等\cite{wallace2024instruction}形式化指令层级（Instruction Hierarchy）概念，训练LLM在面对相互冲突的指令时优先遵循系统指令而非用户输入或外部数据中的指令，通过构建包含指令冲突的训练数据实现30\%以上鲁棒性提升。Chen等\cite{chen2025secalign}提出SecAlign，利用直接偏好优化（DPO）训练模型区分合法指令与注入内容，对每个训练样本构造"正确遵循用户指令"的正例和"误信注入内容"的负例，通过偏好学习使模型倾向正例行为，是首个将GCG攻击成功率降至10\%以下的防御方法。

\textbf{基于检测的防御。}Zhu等\cite{zhu2025melon}提出MELON（Masked Output Evaluation），核心思想是对工具输出进行遮蔽后重新执行智能体——若遮蔽可疑内容后智能体行为发生显著变化，则判定存在注入攻击，在GPT-4o上实现0.24\%攻击成功率，但增加了额外的推理成本。Hung等\cite{hung2025attention}提出AttentionTracker，利用注意力模式中的"分心效应"检测提示注入——当模型对注入内容的注意力权重异常高于正常数据时，表明可能存在注入，该方法无需修改模型即可部署，但对语义隐蔽的注入检测能力有限。Jacob等\cite{jacob2025promptshield}开发PromptShield，训练专门的注入检测分类器，在0.1\%假阳性率下实现65.3\%真阳性率，兼顾了低误报与高检出。Liu等\cite{liu2025datasentinel}提出DataSentinel，以博弈论方法进行极小极大优化训练检测器——将检测器与攻击者建模为零和博弈，检测器最小化最坏情况下的漏检率，在对抗自适应攻击时保持鲁棒。然而，Hackett等\cite{hackett2025bypassing}系统评估了主流护栏系统的鲁棒性，证明许多护栏系统可通过字符注入（如零宽字符、同形字替换）和编码操纵（如Base64、Unicode转义）绕过，揭示基于文本模式匹配的检测方法本质上的脆弱性。

\textbf{架构与系统级防御。}Wang等\cite{wang2025cacheprune}提出CachePrune，基于特征归因的KV缓存神经元剪枝缓解提示注入——识别并剪除对注入内容响应贡献最大的注意力头对应的KV缓存条目，从计算图层面削弱注入对模型行为的影响。Zhong等\cite{zhong2025rtbas}提出RTBAS（Run-Time Based Agent Security），将信息流控制（IFC）应用于LLM智能体，为每个数据源分配机密性和完整性标签，运行时跟踪信息流确保高完整性指令不被低完整性数据污染。Debenedetti等\cite{debenedetti2025defeating}提出从设计上击败提示注入的架构方案CaMeL，通过将智能体操作权限与数据流分离，确保外部数据仅作为信息输入而不能直接触发工具调用。Zverev等\cite{zverev2025can, zverev2025aside}从计算复杂性理论出发，证明LLM在计算上无法可靠区分指令与数据——将该问题归约为不可判定问题，为"完全防御提示注入是不可能的"提供了理论支撑。

\textbf{评估框架。}Debenedetti等\cite{debenedetti2024agentdojo}提出AgentDojo框架，提供跨4个域（工作区、银行、旅行、Slack）的有状态智能体执行环境，包含900+测试用例，通过确定性检查函数（验证工具调用名称和参数是否匹配攻击者目标）评估攻击成功率，而非简单的字符串匹配，是目前最接近真实部署的智能体安全评估平台。Zhan等\cite{zhan2024injecagent}提出InjecAgent基准，评估工具集成LLM智能体中的间接提示注入，包含1054个测试用例，在ReAct提示的GPT-4上报告24\%攻击成功率，但评估场景相对简化。Fang等\cite{fang2025securing}提出FIDES，基于信息流控制的安全AI智能体方案——通过动态污点跟踪追踪机密性和完整性标签的传播，在AgentDojo上验证了该方案可在保证安全性的同时完成广泛的合法任务。

\subsection{现有研究存在的问题}

\begin{enumerate}
  \item \textbf{越狱攻击方法难以直接迁移到提示注入场景。} 现有的自动化对抗攻击方法（如GCG、TAP）主要针对越狱场景设计，其优化目标和成功判定标准与提示注入攻击存在本质差异：越狱关注是否生成被禁止内容（可通过字符串匹配或LLM判断），而提示注入需要触发特定工具调用及正确参数（二元结果，迭代信号有限）。智能体的多轮交互、有状态环境和工具调用决策边界使攻击优化问题更加复杂，直接迁移的攻击成功率显著低于越狱场景。
  \item \textbf{提示注入攻击评估缺乏系统性的自动化方法。} 当前提示注入研究大多依赖手工攻击或简单模板，缺乏在真实智能体环境中进行自动化优化攻击评估的系统性方法。现有评估场景往往限于简化任务（如文本摘要中的恶意指令），而在具有工具调用能力的智能体操作有状态环境中的评估严重不足，无法充分测试智能体对自适应对抗者的鲁棒性。
  \item \textbf{攻击的可迁移性与泛化能力不足。} 在开源小模型上优化的攻击几乎完全无法迁移到先进闭源模型（如GPT-5），成功率低于2\%。跨任务域和跨模型的攻击泛化仍面临根本性挑战：攻击优化的目标与真实部署目标之间存在对齐差距，单一输出序列的优化方式无法覆盖实际中的多种成功输出，且不同模型的工具调用行为和防御特性差异显著。
\end{enumerate}

\subsection{本研究的创新点}

\begin{enumerate}
  \item \textbf{面向LLM智能体间接提示注入的攻击目标形式化与适配。} 将越狱攻击中的攻击目标从"生成被禁止内容"重新形式化为"触发特定工具调用并传递正确参数"的条件概率最大化问题，提出了肯定响应目标（Affirmative Response Target）与LLM翻译目标两种目标字符串构造策略，解决了越狱与提示注入场景间优化目标对齐的核心难题。
  \item \textbf{通用攻击的跨场景联合优化与迁移性评估框架。} 提出跨多任务域的通用攻击联合优化方法——通过在多场景训练集上平均梯度和损失来优化单一注入字符串（或可复用的前缀-后缀组件），并设计保留完全未见域的泛化性评估协议，首次系统性地研究了间接提示注入攻击在智能体场景中的跨任务泛化能力。
  \item \textbf{面向智能体攻击的Tokenization验证与评估可靠性增强机制。} 发现并提出上下文相关分词（context-dependent tokenization）导致的优化-部署编码不一致问题，设计decode-reencode验证过滤器确保token序列的编码-解码往返一致性；针对LLM-as-judge在智能体评估中准确率不足（如Qwen上仅约50\%）的问题，提出多级评估策略与可靠性测试机制，显著提升黑盒攻击的优化信号质量。
\end{enumerate}

\subsection{研究意义与价值}

\subsubsection{理论意义}

\begin{enumerate}
  \item 将越狱攻击中的白盒梯度优化（GCG）与黑盒迭代搜索（TAP）两种主流方法系统性迁移到间接提示注入场景，揭示自动化对抗方法从文本生成到动作执行场景迁移中的理论挑战与性能边界，丰富LLM智能体安全评估的方法论体系。
  \item 构建通用攻击优化框架，研究在多任务、多场景联合优化下攻击的泛化机制，探索通用注入模式的存在性与可发现性，为理解LLM智能体的系统性脆弱点提供理论支撑。
  \item 通过全面的消融实验与迁移性分析，揭示影响攻击成功的关键因素（注入结构、语义内容、初始化随机性、评估器可靠性等），为智能体安全防御的设计提供攻击侧的理论依据。
\end{enumerate}

\subsubsection{实用价值}

\begin{enumerate}
  \item 扩展AgentDojo框架以支持自动化攻击优化，集成HuggingFace Transformers库实现白盒梯度攻击，开发标准化攻击集成接口，为LLM智能体安全评估提供可直接使用的开源工具。
  \item 在跨多域（工作区、银行、旅行、Slack）和多模型（Qwen3-4B、Gemma-3-4B、GPT-5）的80个任务对上进行全面实证评估，为工业界部署LLM智能体提供量化的安全风险评估参考。
  \item 针对开源模型高达45.2\%攻击成功率和72.5\%安全阈值突破率的实验结果，为安全工程师设计防御策略、评估护栏系统有效性提供实证依据和攻击基线。
\end{enumerate}

\subsubsection{与国家及行业发展战略的契合度}
\begin{enumerate}
  \item 契合《新一代人工智能发展规划》中关于"人工智能安全"的战略部署，为AI系统的安全评估与风险防控提供关键技术支撑。
  \item 响应《网络安全法》和《数据安全法》对关键信息基础设施安全保护的要求，为LLM智能体在金融、政务等敏感领域的安全部署提供攻击评估方法。
  \item 服务于国家人工智能标准化总体组关于"人工智能可信特征"标准的研究，为LLM智能体鲁棒性与安全性评测提供方法论参考。
\end{enumerate}


